\scanpage{25}
\uline{1.1} Soit $n\in\mathbb N^*$, et posons
$C_n\subset\mathbb P^2_{\mathbb Q}$ équation projective
\[
(1)\qquad x^n+y^n+z^n=0
\]
Le groupe $\mu_{n\mathbb Q}^3$ opère sur $\mathbb E^3_{\mathbb Q}$ par matrices diagonales,
donc posant
\[
(2)\qquad \Pi_{n\mathbb Q}=\mu_n^3/\underset{\substack{\curvearrowright\,\text{diag.}}}{\mu_n}
\simeq\mu_n\otimes_{\mathbb Z}(\mathbb Z^3/\mathbb Z)
\]
\margin{NB on écrit $\mu_n$ au lieu de $\mu_{n\mathbb Q}$}
opère sur $\mathbb P^2_{\mathbb Q}$, en laissant $C_n$ invariant. La courbe
$C_n/\Pi_{n\mathbb Q}$ s'identifie à la courbe \inserted{$C_1$} qui est donnée,
en coordonnées homogènes, par
\[
(3)\qquad X+Y+Z=0
\]
\margin{Comme carte pondérée, elle est formée de $2n^2$ triangles
ch. sommet est d'ordre $2n$, chaque type de sommet \unsure{intervient} $n$ fois\ldots}
qu'on identifie à $\mathbb P^1_{\unsure{x}}$, en envoyant les points de $C_1$
pour lesquels $X,Y$ ou $Z$ sont nuls resp. en $0,1,\infty$.
Donc écrivant $C_1$ en coordonnées affines, faisant $Z=1$
i.e. $X/Z=\xi$, $Y/Z=\eta$, sous la forme
\[
(3')\qquad \xi+\eta+1=0
\]
i.e.
\[
\eta=-1-\xi
\]
$\xi=0$ correspond à $\eta=-1$, $\eta=0$ à $\xi=-1$, donc posant
\[
(3'')\qquad t=-\xi\qquad\text{donc}\quad \xi=-t,\ \eta=t-1,\ \zeta=1
\]
les valeurs de $t$ pour lesquelles $X=0$, $Y=0$, $Z=0$ sont resp. $t=0,1,\infty$.

\margin{Il y a en $\unsure{\alpha(i)|\alpha_i}$ $(i=0,1,\infty)$
$n$ p\textsuperscript{ts} de ramification de $C_n$, chacun de degré de ramification $n$.}
$C_n$ est étale sur $C_1=\mathbb P^1_{\mathbb Q}$ sauf en $0,1,\infty$, \unsure{un|on}
\unsure{donne|trouve} rev. principal $C_n^*$ de $C_1^*\simeq\mathbb P^1_{\mathbb Q}\setminus\{0,1,\infty\}$,
groupe $\Pi_{n\mathbb Q}$.
% archive foliation: 24
